\section*{Capítulo \ref{ch:b3:behavioural-ecology} --- Ecologia comportamental}

\begin{solution}{exo:b3:behavioural-ecology:1}
Mecanismo: a dançarina transpõe o ângulo entre a comida e o sol, medido
durante o voo, para o ângulo em relação à gravidade no favo, e codifica a
distância a partir do fluxo óptico da jornada como a duração da corrida.
Desenvolvimento: em grande parte inato --- abelhas criadas sem ver danças
dançam corretamente, com pequenos ajustes aprendidos. Função: recruta
companheiras de ninho para uma fonte lucrativa, elevando a entrada de
alimento da colônia. História: derivada de movimentos de intenção
ritualizados de partir para a comida; as abelhas-anãs ainda dançam numa
superfície horizontal apontando diretamente para a fonte, a forma
ancestral.
\end{solution}

\begin{solution}{exo:b3:behavioural-ecology:2}
Um \omterm{def:b3:behavioural-ecology:tinbergen}{padrão fixo de ação} é uma sequência estereotipada disparada por uma
pista simples, o \omterm{def:b3:behavioural-ecology:tinbergen}{estímulo-sinal}, e completada uma vez iniciada: o rolar
do ovo do ganso, disparado por um objeto em forma de ovo; a bicada do
filhote de gaivota, disparada por uma mancha vermelha numa forma
parecida com um bico. O bastão com três faixas vermelhas, bicado mais que
uma cabeça de verdade, mostrou que o mecanismo disparador é um filtro
sintonizado numa característica --- contraste vermelho, alongamento ---,
e não um reconhecimento do pai, e que exagerar a característica exagera a
resposta.
\end{solution}

\begin{solution}{exo:b3:behavioural-ecology:3}
Deixe uma mancha quando a taxa de ganho nela tiver caído à taxa média do
hábitat como um todo, viagem incluída. Dobrar o tempo de viagem baixa
essa média, de modo que o forrageador fica mais tempo em cada mancha,
esgota cada uma mais a fundo e termina com uma taxa global menor.
\end{solution}

\begin{solution}{exo:b3:behavioural-ecology:4}
Uma EEE é uma estratégia que, uma vez que quase todos a jogam, nenhuma
alternativa rara consegue superar. Só-Pomba não é uma: um Falcão
solitário vence toda disputa e ganha $V$ contra o $V/2$ das Pombas. Na
mistura cada estratégia ganha $V(1 - V/C)/2$, menos que o $V/2$ que
só-Pomba daria, porque uma fração das disputas são lutas que ferem;
ninguém consegue melhorar isso sozinho, de modo que todos ficam com o
prejuízo.
\end{solution}

\begin{solution}{exo:b3:behavioural-ecology:5}
$V = 6$, $C = 10$: $p^{*} = 0.6$; ganho $6(1 - 0.6)/2 = 1.2$ para cada
uma, contra $3$ em só-Pomba. $V = 12 > C$: Falcão é a EEE pura, com ganho
$(12 - 10)/2 = 1$, contra $6$ em só-Pomba --- a disputa destrói cinco
sextos do valor.
\end{solution}

\begin{solution}{exo:b3:behavioural-ecology:6}
$\tau = T_{0}$: $t^{*} = 1.15\times 2 = \qty{2.3}{min}$, fração $1 -
\mathrm{e}^{-1.15} = 0.68$, $R = 0.68A/4.3 = 0.16A$ por minuto. $\tau =
4T_{0}$: $t^{*} = 2.0\times 2 = \qty{4}{min}$ (mais precisamente $3.9$),
fração $0.86$, $R = 0.86A/12 = 0.072A$ por minuto.
\end{solution}

\begin{solution}{exo:b3:behavioural-ecology:7}
Meios-irmãos, $1/4$; avô--neto, $1/4$; primos de primeiro grau, $1/8$.
Uma operária haplodiploide e seu irmão: ele é haploide e só tem os genes
da mãe dos dois; dos genes da operária, metade veio da mãe e metade
desses são as cópias que ele recebeu, de modo que $r = 1/4$ do ponto de
vista dela (do dele é $1/2$: o parentesco não é simétrico quando as
ploidias diferem). Rainha e filho: todos os genes dele são dela, mas só
metade dos dela estão nele: $r = 1/2$ do lado dela, $1$ do dele.
\end{solution}

\begin{solution}{exo:b3:behavioural-ecology:8}
$n\,r\,b > c$ com $b = 0.01$ e $c = 0.02$: $n\,r > 2$. Irmãos completos,
$r
= 1/2$: $n \ge 5$. Sobrinhos, $r = 1/4$: $n \ge 9$. Não aparentados:
nunca --- e é por isso que os machos que se dispersam não chamam.
\end{solution}

\begin{solution}{exo:b3:behavioural-ecology:9}
Os acasalamentos dobram a cada \qty{25}{cm} e a sobrevivência cai um
quinto: $L = 25$: $0.5\times 1.2 = 0.6$; $50$: $1$; $75$: $2\times 0.8 = 1.6$; $100$: $4\times 0.6 = 2.4$; $125$: $8\times 0.4 = 3.2$; $150$:
$16\times 0.2 =
3.2$. Nesses termos o produto continua subindo muito além
dos \qty{50}{cm} naturais, de modo que o modelo está deixando custos de
fora: a penalidade de sobrevivência de uma cauda longa em voo é mais que
linear, a resposta feminina satura, fazer crescer as penas custa energia,
e o dobro de ninhos de Andersson foi medido ao longo de uma estação, e
não de uma vida. A cauda natural está onde os custos reais mordem.
\end{solution}

\begin{solution}{exo:b3:behavioural-ecology:10}
Retaliadora: contra Falcão, $(V - C)/2$; contra Pomba, $V/2$; contra
Retaliadora, $V/2$. Numa população só de Retaliadoras, um Falcão raro
encontra Retaliadoras, que revidam, e ganha $(V - C)/2 < V/2$: ele não
consegue invadir quando $V < C$. Uma Pomba rara ganha $V/2$, exatamente o
ganho das Retaliadoras, de modo que as Pombas são neutras e podem entrar
por deriva; uma vez comuns, elas deixam os Falcões entrar (um Falcão
ganha $V$ contra uma Pomba), e o jogo de três estratégias não tem EEE
simples --- um primeiro exemplo de como acrescentar uma opção pode
desestabilizar uma mistura estável.
\end{solution}

\begin{solution}{exo:b3:behavioural-ecology:11}
Se todos persistissem exatamente $t_{0}$, um mutante que persistisse
$t_{0} +
\varepsilon$ venceria toda disputa por um custo extra
desprezível, e uma população desses mutantes seria batida por $t_{0} + 2\varepsilon$, e assim por diante: nenhum tempo fixo é estável, de modo
que a EEE precisa ser um tempo aleatório. Com uma distribuição
exponencial, a probabilidade de persistir mais $s$, dada uma persistência
de $t$ até agora, é $\mathrm{e}^{-s/m}$ qualquer que seja $t$ --- a
propriedade de falta de memória ---, de modo que um oponente ainda em
exibição nada revela sobre quanto tempo mais vai continuar, e nenhuma
estratégia consegue explorar o tempo decorrido.
\end{solution}

\begin{solution}{exo:b3:behavioural-ecology:12}
As operárias valorizam uma irmã em $3/4$ e um irmão em $1/4$, de modo que
sob controle das operárias a colônia deveria investir três vezes mais em
fêmeas que em machos, uma razão de $3{:}1$ por investimento; a rainha,
igualmente aparentada aos dois, prefere $1{:}1$. Trivers e Hare (1976)
encontraram razões de investimento perto de $3{:}1$ em formigas com uma
única rainha de acasalamento único: as operárias vencem. O acasalamento
múltiplo baixa o parentesco médio de uma operária com suas irmãs para
algo entre $1/4$ e $1/2$ e empurra o ótimo das operárias rumo a $1{:}1$;
as colônias de espécies com rainhas de muitos acasalamentos investem, de
fato, de modo mais equilibrado, como o argumento prevê.
\end{solution}

\begin{solution}{pb:b3:behavioural-ecology:1}
\textbf{1.} $R(t) = A(1 - \mathrm{e}^{-t/T_{0}})/(\tau + t)$; sair quando
$g'(t) = R(t)$: $(A/T_{0})\mathrm{e}^{-t/T_{0}} = A(1 - \mathrm{e}^{-t/T_{0}})/
(\tau + t)$.
\textbf{2.} Multiplique por $(\tau + t)\mathrm{e}^{t/T_{0}}/A$: $(\tau + t)/T_{0}
= \mathrm{e}^{t/T_{0}} - 1$. Com $\tau = T_{0}$ e $x = t/T_{0}$: $2 + x =
\mathrm{e}^{x}$; $x = 1.15$ dá $3.15$ contra $3.16$.
\textbf{3.} $\tau = 4T_{0}$: $5 + x = \mathrm{e}^{x}$; $x = 1.94$ dá
$6.94$ contra $6.96$; $t^{*} \approx 1.94\,T_{0}$, digamos $2.0$.
\textbf{4.} Denso: $t^{*} = \qty{3.5}{min}$, $60(1 - \mathrm{e}^{-1.15}) =
\qty{41}{kJ}$, \qty{68}{\%}, $R = 41/6.5 = \qty{6.4}{kJ/min}$. Esparso:
$t^{*} = \qty{5.8}{min}$, \qty{51}{kJ}, \qty{86}{\%}, $R = 51/17.8 =
\qty{2.9}{kJ/min}$.
\textbf{5.} Cinco minutos: $g(5) = 60(1 - \mathrm{e}^{-5/3}) = \qty{49}{kJ}$.
Denso: $49/8 = \qty{6.1}{kJ/min}$ (\qty{96}{\%} do ótimo); esparso:
$49/17 = \qty{2.9}{kJ/min}$ (\qty{99}{\%}). Uma regra grosseira perde
alguns por cento --- o bastante para a seleção ter se contentado com ela.
\textbf{6.} O intervalo entre capturas se alonga à medida que a mancha se
esgota, de modo que um tempo de desistência fixo é uma taxa marginal de
captura fixa --- o critério do teorema, que é o mesmo em toda mancha. Ele
falha porque o limiar certo depende do hábitat: com \qty{20}{s} fixos o
animal sai cedo demais onde a viagem é longa e tarde demais onde ela é
curta, a menos que o próprio tempo de desistência seja ajustado à
experiência recente --- o que os forrageadores fazem.
\textbf{7.} Falcão contra Falcão, $(10 - 30)/2 = -10$; Falcão contra
Pomba, $10$; Pomba contra Falcão, $0$; Pomba contra Pomba, $5$; $p^{*} = V/C = 1/3$.
\textbf{8.} $W_{H} = 10 - 20p$, $W_{D} = 5(1 - p)$. $p = 0.1$: $8$ contra
$4.5$, os Falcões se espalham. $p = 1/3$: $3.33$ cada, equilíbrio. $p = 0.6$: $-2$ contra $2$, as Pombas se espalham.
\textbf{9.} $3.33$ na EEE contra $5$ em só-Pomba: um terço do valor de
cada disputa se perde pela possibilidade de lutar.
\textbf{10.} $V = 40 > C$: Falcão é a EEE pura; toda disputa é uma luta.
Os anos difíceis deveriam ver mais lutas e mais lesões --- como de fato
veem.
\textbf{11.} Um Falcão que encontra uma Burguesa a acha residente metade
das vezes (uma luta, $(V - C)/2$) e invasora metade das vezes (uma vitória
fácil, $V$): média $(3V - C)/4$. A Burguesa ganha $V/2$ entre as suas. O
Falcão só invade se $(3V - C)/4 > V/2$, isto é, se $V > C$. O ``o
residente vence'' da borboleta é Burguesa: uma convenção barata, estável
porque lutar custa mais do que vale uma mancha de sol --- e, quando os
dois se creem residentes, os dois jogam Falcão.
\textbf{12.} A melhor coisa a fazer depende do que todos os outros fazem,
de modo que a população se acomoda onde as estratégias pagam igual, e não
no melhor de ninguém.
\textbf{13.} $n\bar r\,b > c$: $n\bar r > c/b = 3$.
\textbf{14.} Oito irmãos completos: $n\bar r = 4 > 3$, o turno compensa.
Oito estranhos: $0$, nunca. Oito meios-irmãos: $2 < 3$, não.
\textbf{15.} $c = 0.01$: limiar $n\bar r > 1$ --- meios-irmãos e até
alguns irmãos completos já bastam. O indivíduo de menor custo é aquele
para quem a regra se satisfaz, de modo que os bem alimentados montam
guarda enquanto os famintos cavam.
\textbf{16.} $k$ irmãs valem $\tfrac{3}{4}k$, e $k$ filhas,
$\tfrac{1}{2}k$: prefira irmãs, na razão $3{:}2$.
\textbf{17.} Uma irmã ao acaso é irmã completa ($3/4$) com probabilidade
$1/3$ e meia-irmã ($1/4$) com probabilidade $2/3$: $\bar r = 0.25 +
0.17 = 0.42 < 0.5$. A preferência se inverte: as filhas agora batem as irmãs,
de modo que o parentesco sozinho não explica as operárias estéreis nas
espécies poliândricas --- é preciso acrescentar a ecologia e o
policiamento no nível da colônia.
\textbf{18.} A regra só exige que $r$ seja em média maior entre os
ajudados que na população; qualquer pista correlacionada com o parentesco
serve --- ficar perto da toca natal, ajudar aqueles com quem se cresceu.
A regra de um esquilo-terrestre pode ser ``chame quando estiver perto de
casa'', o que separa os parentes sem reconhecer nenhum.
\textbf{19.} $W = (L/20)(1 - L^{2}/10^{4})$: $L = 20$: $0.96$; $40$:
$1.68$; $60$: $1.92$; $80$: $1.44$.
\textbf{20.} $W' = \tfrac{1}{20}(1 - 3L^{2}/10^{4}) = 0$: $L^{*} = 100/\sqrt{3}
= \qty{58}{cm}$, $W = 2.89\times\tfrac{2}{3} = 1.92$.
\textbf{21.} $L^{*} = 70/\sqrt{3} = \qty{40}{cm}$: mais curta. O ótimo
fica onde o ganho marginal em acasalamentos iguala a perda marginal em
sobrevivência, de modo que um custo de predação mais pesado encolhe a
cauda, e populações sob predações diferentes deveriam diferir no ornamento
--- como diferem os guppies e os espadas.
\textbf{22.} As duas. A \omterm{def:b3:behavioural-ecology:sexual-selection}{seleção desenfreada} prevê que a preferência
ultrapassa a característica, que para no ótimo de sobrevivência; o
handicap prevê que as fêmeas preferem mais longa porque só um macho em
forma a carrega, de novo além do que a maioria consegue fazer crescer. O
experimento mostra uma preferência além do ótimo e não decide entre as
explicações.
\textbf{23.} Se a cauda fosse gratuita, todo macho poderia fazer crescer
a mais longa, ela nada diria sobre sua qualidade, e uma preferência por
ela seria explorada e decairia. Um custo que pesa mais sobre um macho em
má condição torna a cauda longa acessível apenas aos em forma, e a
preferência então seleciona aptidão: o custo é o que torna o \omterm{prop:b3:behavioural-ecology:signals}{sinal
honesto}.
\textbf{24.} Com oito vezes o retorno por acasalamento, a aptidão de um
macho sobe acentuadamente com os acasalamentos e a de uma fêmea quase
nada; os machos são selecionados para competir por acasalamentos --- as
disputas Falcão--Pomba e suas armas --- e as fêmeas, cujos poucos
descendentes contam cada um, são selecionadas para escolher, o que faz os
ornamentos masculinos compensarem. Uma assimetria no ganho por
acasalamento produz tanto a briga quanto o enfeite.
\textbf{25.} $t^{*} \approx \qty{3.5}{min}$ no hábitat denso; $p^{*}
= 1/3$; o turno de sentinela compensa quando $n\bar r > 3$.
\end{solution}
